\documentclass{article}
\usepackage[T1]{fontenc}
\usepackage{lmodern}
\usepackage{graphicx}
\usepackage{amsmath,amssymb}
\usepackage[a4paper,margin=2cm]{geometry}
\usepackage[absolute,overlay]{textpos}
\setlength{\TPHorizModule}{1mm}
\setlength{\TPVertModule}{1mm}
\usepackage[indonesian]{babel}
\usepackage{hyperref}
\setlength{\emergencystretch}{2em}
% unit: Halaman 23

\begin{document}

% === Halaman 23 (llm) ===

Andaikan kriptanalis mengetahui ukuran blok = 2 karakter (16 bit), spasi diabaikan.

Blok-blok cipherteks:
\[ C_{1}, C_{2}, C_{3}, C_{4}, C_{5}, C_{6}, C_{7}, \color{red}{C_{8}, C_{9}}, C_{10}, C_{11}, C_{12}, C_{13}, C_{14}, C_{15}, C_{16} \]

Misalkan kriptanalis mengethaui posisi pesan yang berkaitan dengan nominal uang, yaitu blok $C_{8}$ dan $C_{9}$.

Kriptanalis membuang blok cipherteks $C_{8}$ dan $C_{9}$ :
\[ C_{1}, C_{2}, C_{3}, C_{4}, C_{5}, C_{6}, C_{7}, C_{10}, C_{11}, C_{12}, C_{13}, C_{14}, C_{15}, C_{16} \]

\end{document}
